\documentclass[12pt]{article}

\usepackage[utf8]{inputenc}
\usepackage[T1]{fontenc}
\usepackage{lmodern}
\usepackage[spanish]{babel}
\usepackage{booktabs}
\usepackage{amsmath}
\usepackage{forest}
\usepackage{float}
\usepackage{listings}
\usepackage{xcolor}
\usepackage{tikz}

\definecolor{codegreen}{rgb}{0,0.6,0}
\definecolor{codegray}{rgb}{0.5,0.5,0.5}
\definecolor{codepurple}{rgb}{0.58,0,0.82}
\definecolor{backcolour}{rgb}{0.95,0.95,0.92}

\lstdefinestyle{mystyle}{
    backgroundcolor=\color{backcolour},   
    commentstyle=\color{codegreen},
    keywordstyle=\color{magenta},
    numberstyle=\tiny\color{codegray},
    stringstyle=\color{codepurple},
    basicstyle=\ttfamily\footnotesize,
    breakatwhitespace=false,         
    breaklines=true,                 
    captionpos=b,                    
    keepspaces=true,                 
    numbers=left,                    
    numbersep=5pt,                  
    showspaces=false,                
    showstringspaces=false,
    showtabs=false,                  
    tabsize=2
}

\lstset{style=mystyle}

\sloppy
\setlength{\parindent}{0pt}

\begin{document}

% Título y materia
\begin{center}
  {\LARGE \textbf{Random Forest}}\\[0.5em]
  {\normalsize (Bosques Aleatorios)}\\[0.5em]
  {Analítica de Datos, Universidad de San Andrés}
\end{center}

Si encuentran algún error en el documento o hay alguna duda, mandenmé un mail a rodriguezf@udesa.edu.ar y lo revisamos.

\section{¿Qué es Random Forest?}
Random Forest es básicamente un montón de árboles de decisión trabajando juntos. ¿Por qué? Porque un solo árbol puede ser muy caprichoso y ajustarse demasiado a los datos de entrenamiento (overfitting). Es como si un solo árbol dijera ``bueno, vi estos datos y ahora me sé todo'', pero cuando llegan datos nuevos, se confunde. La idea de Random Forest es simple: en lugar de confiar en un solo árbol, construimos muchos árboles y que voten. Es como preguntarle a un montón de amigos su opinión en lugar de solo a uno. La mayoría probablemente tenga razón.

\vspace{1em}

\textit{¿Por qué funciona mejor?} Porque cada árbol ve los datos de manera diferente (usando subconjuntos aleatorios), entonces cuando se equivocan, se equivocan de maneras distintas. Al promediar sus respuestas, los errores se cancelan y obtenemos una predicción más estable.

\section{¿Cómo funciona el algoritmo?}
El algoritmo de Random Forest es bastante directo:

\begin{enumerate}
    \item \textbf{Bootstrap:} Para cada árbol, tomamos una muestra aleatoria de nuestros datos (con reemplazo). Esto significa que algunos datos podrían aparecer varias veces y otros no aparecer en absoluto.
    \item \textbf{Selección aleatoria de variables:} En cada división del árbol, solo consideramos un subconjunto aleatorio de las variables disponibles. Esto evita que una variable muy importante domine todos los árboles.
    \item \textbf{Construcción del árbol:} Construimos el árbol completo usando CART con los datos y variables seleccionadas.
    \item \textbf{Repetir:} Hacemos esto muchas veces (típicamente 100-1000 árboles).
    \item \textbf{Votación:} Para predecir, cada árbol vota y tomamos la decisión de la mayoría (clasificación) o el promedio (regresión).
\end{enumerate}

\section{Ventajas de Random Forest}
Random Forest tiene varias ventajas que lo hacen muy popular:

\begin{itemize}
    \item \textbf{Reduce el overfitting:} Al promediar muchos árboles, suavizamos las decisiones y evitamos ajustarnos demasiado a los datos de entrenamiento.
    \item \textbf{Maneja bien diferentes tipos de variables:} No le importa si son numéricas o categóricas, las maneja todas bien.
    \item \textbf{Es robusto:} Los outliers y el ruido en los datos no lo afectan tanto como a otros algoritmos.
    \item \textbf{Importancia de variables:} Nos dice qué variables son más importantes para hacer predicciones.
    \item \textbf{No necesita preprocesamiento:} No tenemos que normalizar los datos ni hacer transformaciones complicadas.
    \item \textbf{Es rápido:} Aunque construye muchos árboles, cada uno es simple y se puede hacer en paralelo.
\end{itemize}

\section{Parámetros importantes}
Cuando usamos Random Forest, hay algunos parámetros que podemos ajustar:

\begin{itemize}
    \item \textbf{n\_estimators:} Cuántos árboles queremos (más árboles = mejor rendimiento pero más lento). Típicamente entre 100-1000.
    \item \textbf{max\_features:} Cuántas variables considerar en cada división. Por defecto es $\sqrt{\text{número de variables}}$ para clasificación.
    \item \textbf{max\_depth:} Profundidad máxima de cada árbol. Si no la ponemos, los árboles crecen hasta que no se pueden dividir más.
    \item \textbf{min\_samples\_split:} Mínimo de muestras necesarias para dividir un nodo.
    \item \textbf{min\_samples\_leaf:} Mínimo de muestras en cada hoja final.
    \item \textbf{bootstrap:} Si queremos usar bootstrap (por defecto es True).
\end{itemize}

\textit{¿Cómo elegimos estos parámetros?} La mayoría de las veces, los valores por defecto funcionan bien. Si queremos optimizar, podemos usar validación cruzada para encontrar la mejor combinación.

\section{Implementación en Python}
\begin{lstlisting}[language=Python]
# Importamos las librerias necesarias
from sklearn.ensemble import RandomForestClassifier
from sklearn.model_selection import train_test_split
from sklearn.metrics import accuracy_score, classification_report
from sklearn.datasets import load_wine

# Cargamos el dataset Wine
wine = load_wine()
X = wine.data
y = wine.target

# Dividimos en entrenamiento y prueba
X_train, X_test, y_train, y_test = train_test_split(
    X, y, test_size=0.2, random_state=42
)

# Creamos el modelo de Random Forest
rf_model = RandomForestClassifier(
    n_estimators=100,        # 100 arboles
    max_features='sqrt',     # sqrt del numero de variables
    random_state=42
)

# Entrenamos el modelo
rf_model.fit(X_train, y_train)

# Hacemos predicciones
y_pred = rf_model.predict(X_test)

# Evaluamos el rendimiento
accuracy = accuracy_score(y_test, y_pred)
print(f'Precision: {accuracy:.3f}')

# Vemos el reporte detallado
print(classification_report(y_test, y_pred))
\end{lstlisting}

\section{Importancia de variables}
Una de las cosas más útiles de Random Forest es que nos dice qué variables son más importantes:

\begin{lstlisting}[language=Python]
# Obtenemos la importancia de las variables
importances = rf_model.feature_importances_
feature_names = wine.feature_names

# Creamos un DataFrame para ver mejor
import pandas as pd
importance_df = pd.DataFrame({
    'Variable': feature_names,
    'Importancia': importances
}).sort_values('Importancia', ascending=False)

print(importance_df)

# Graficamos la importancia
import matplotlib.pyplot as plt
plt.figure(figsize=(10, 6))
plt.barh(importance_df['Variable'], importance_df['Importancia'])
plt.xlabel('Importancia')
plt.title('Importancia de Variables en Random Forest')
plt.show()
\end{lstlisting}

\section{Comparación con CART}
¿Por qué Random Forest es mejor que un solo árbol CART?

\begin{table}[H]
    \centering
    \begin{tabular}{p{4cm}|p{4cm}|p{4cm}}
        \toprule
        \textbf{Característica} & \textbf{CART} & \textbf{Random Forest} \\
        \midrule
        Overfitting & Propenso & Menos propenso \\
        Estabilidad & Baja & Alta \\
        Interpretabilidad & Alta & Media \\
        Velocidad & Rápido & Más lento \\
        Precisión & Variable & Generalmente mejor \\
        \bottomrule
    \end{tabular}
\end{table}

\textit{¿Cuándo usar cada uno?} Si necesitamos entender exactamente cómo toma las decisiones, CART es mejor. Si queremos la mejor precisión posible y no nos importa tanto la interpretabilidad, Random Forest es la opción.

\section{Notebook Anexo}
En el repositorio hay un notebook anexo en el que usamos el dataset \textbf{Wine} de scikit-learn, que contiene información química de 178 vinos de tres variedades diferentes. Cada observación tiene 13 variables numéricas (alcohol, magnesio, fenoles, color, etc.) y una variable objetivo que indica la clase de vino (0, 1 o 2).

\vspace{0.5em}
Para detalles completos de las variables y análisis exploratorio, consultar el notebook. En el experimento:

\begin{itemize}
    \item \textbf{CART:} Se equivocó en 2 casos de 36 (94.4\% de precisión)
    \item \textbf{Random Forest:} Se equivocó en 0 casos de 36 (100\% de precisión)
\end{itemize}

\textit{¿Por qué Random Forest fue mejor?} Porque al promediar las decisiones de 100 árboles, suavizamos los errores individuales y obtenemos una predicción más robusta. Es como tener 100 expertos en vino en lugar de uno solo.

\vspace{1em}

\textit{¿Qué pasa si aumentamos el test size?} Probablemente Random Forest siga siendo mejor, pero la diferencia podría ser menor. La clave está en que Random Forest generaliza mejor a datos nuevos.

\end{document}
